\documentclass[11pt]{article}

\usepackage[margin=1in]{geometry}
\usepackage{amsmath, amssymb}
\usepackage{xcolor}
\usepackage{listings}
\usepackage{hyperref}
\usepackage{array}
\usepackage{booktabs}

\hypersetup{
    colorlinks=true,
    linkcolor=blue!60!black,
    urlcolor=blue!60!black,
    citecolor=blue!60!black
}

\lstdefinestyle{bugpython}{
  language=Python,
  basicstyle=\ttfamily\small,
  keywordstyle=\color{blue!60!black},
  stringstyle=\color{green!40!black},
  commentstyle=\color{gray!70!black},
  showstringspaces=false,
  breaklines=true,
  frame=single,
  rulecolor=\color{gray!30},
  columns=fullflexible,
  keepspaces=true
}

\title{SymPy Bug Report}
\author{}
\date{}

\begin{document}

\maketitle

\section{Bug 60: LerchPhi Left Limit Returns Divergent Endpoint}

\subsection{Minimal Reproducer}

\medskip

\begin{lstlisting}[style=bugpython]
from sympy import lerchphi, limit, symbols

x = symbols("x")
print(limit(lerchphi(x, 1, 1), x, 1, dir="-"))
\end{lstlisting}

\subsection{Observed Behavior}

\medskip

\begin{lstlisting}[style=bugpython]
limit = lerchphi(1, 1, 1)
\end{lstlisting}

On the checked SymPy development version, the left-hand limit is returned as the endpoint expression \(\operatorname{lerchphi}(1,1,1)\). That endpoint is not a finite value of the original analytic expression; it is the divergent defining series after substituting \(x=1\).

\subsection{Expected Behavior}

The correct result is \(+\infty\):
\[
\lim_{x\to 1^-} \operatorname{lerchphi}(x,1,1)=\infty .
\]
For the representative case, applying SymPy's functional expansion first gives
\[
\operatorname{lerchphi}(x,1,1)=\frac{-\log(1-x)}{x},
\]
whose left-hand limit at \(x=1\) is \(+\infty\).

\subsection{Mathematical Explanation}

For \(0<x<1\), the Lerch transcendent has the convergent series representation
\[
\operatorname{lerchphi}(x,1,a)=\sum_{n=0}^{\infty}\frac{x^n}{n+a}.
\]
In the representative case \(a=1\),
\[
\operatorname{lerchphi}(x,1,1)
=\sum_{n=0}^{\infty}\frac{x^n}{n+1}
=\frac{-\log(1-x)}{x}.
\]
As \(x\to 1^-\), the numerator \(-\log(1-x)\) diverges to \(+\infty\) and the denominator tends to \(1\). Therefore the left-hand limit is \(+\infty\), not an unevaluated endpoint.

The returned expression \(\operatorname{lerchphi}(1,1,1)\) is not an equivalent representation of the limit. At \(x=1\) the defining series becomes
\[
\sum_{n=0}^{\infty}\frac{1}{n+1},
\]
the harmonic series, which diverges. Returning the endpoint expression masks that divergence instead of reporting the limit.

\subsection{Source-Level Diagnosis}

The source-level diagnosis locates the immediate wrong step in \nolinkurl{sympy/series/gruntz.py}, specifically in \texttt{mrv\_leadterm} around lines 520--538. The call path described in the artifact is: \texttt{limit(...)} constructs a \texttt{Limit} object and enters \texttt{Limit.doit} in \nolinkurl{sympy/series/limits.py}; after the ordinary leading-term path cannot expand the shifted expression, \texttt{Limit.doit} calls Gruntz. For the left-hand finite limit, Gruntz substitutes \(x=1-1/z\) and analyzes
\[
\operatorname{lerchphi}\left(1-\frac{1}{z},1,1\right)
\]
as \(z\to\infty\).

The diagnosis identifies the following traced behavior:

\begin{lstlisting}[style=bugpython]
e0 = lerchphi(1 - 1/z, 1, 1)
e0.rewrite("tractable", deep=True, limitvar=z) -> e0
mrv_leadterm(e0, z) -> (lerchphi(1, 1, 1), 0)
limitinf(e0, z) -> lerchphi(1, 1, 1)
\end{lstlisting}

Thus the generic Gruntz leading-term machinery treats the special function as having exponent \(0\) and coefficient \(\operatorname{lerchphi}(1,1,1)\). The subsequent \texttt{limitinf} logic returns that coefficient unchanged, so the endpoint is substituted at a divergence point.

\begin{sloppypar}
The special-function implementation contributes to the miss because \nolinkurl{sympy/functions/special/zeta_functions.py} contains expansion rules that can expose this case: \(\operatorname{lerchphi}(x,1,1)\) expands to \(-\log(1-x)/x\). However, the diagnosis says \texttt{lerchphi} does not provide an \texttt{eval}, leading-term, or tractable rewrite hook that this limit path applies automatically near \(x=1\). A plausible fix direction identified by the artifact is to add limit, leading-term, or tractable rewrite support for \(\operatorname{lerchphi}\) near \(z=1\), at least for the \(a=1\) specialization and the \(s=1\) logarithmic case, so the generic Gruntz path does not treat the divergent endpoint as a finite coefficient.
\end{sloppypar}

\subsection{Additional Failing Instances}

The independent verification checks the same left-hand limit for \(a=1,2,3,4,5,6\). In each case, SymPy returns the endpoint expression \(\operatorname{lerchphi}(1,1,a)\), while the expected result is \(+\infty\).

\begin{center}
\renewcommand{\arraystretch}{1.3}
\begin{tabular}{@{}>{\raggedright\arraybackslash}p{0.41\linewidth}>{\raggedright\arraybackslash}p{0.26\linewidth}>{\raggedright\arraybackslash}p{0.23\linewidth}@{}}
\toprule
\textbf{Input / Expression} & \textbf{SymPy behavior} & \textbf{Expected behavior} \\
\midrule
\(\mathrm{limit}(\operatorname{lerchphi}(x,1,1), x, 1, -)\) & \(\operatorname{lerchphi}(1,1,1)\) & \(+\infty\) \\
\(\mathrm{limit}(\operatorname{lerchphi}(x,1,2), x, 1, -)\) & \(\operatorname{lerchphi}(1,1,2)\) & \(+\infty\) \\
\(\mathrm{limit}(\operatorname{lerchphi}(x,1,3), x, 1, -)\) & \(\operatorname{lerchphi}(1,1,3)\) & \(+\infty\) \\
\(\mathrm{limit}(\operatorname{lerchphi}(x,1,4), x, 1, -)\) & \(\operatorname{lerchphi}(1,1,4)\) & \(+\infty\) \\
\(\mathrm{limit}(\operatorname{lerchphi}(x,1,5), x, 1, -)\) & \(\operatorname{lerchphi}(1,1,5)\) & \(+\infty\) \\
\(\mathrm{limit}(\operatorname{lerchphi}(x,1,6), x, 1, -)\) & \(\operatorname{lerchphi}(1,1,6)\) & \(+\infty\) \\
\bottomrule
\end{tabular}
\end{center}

The expected behavior is the same across this family. For every listed positive integer \(a\), the terms are positive on \(0<x<1\), and as \(x\to 1^-\) the series tends to the divergent shifted harmonic series
\[
\sum_{n=0}^{\infty}\frac{1}{n+a}.
\]
The verification script also records large finite partial sums near \(x=1\), consistent with monotone growth toward \(+\infty\), not convergence to a finite endpoint value.

\subsection{Related Issues Assessment}

The deduplication report classifies this as a new instance of a known failure mode. The closest public matches involve LerchPhi, summation endpoint behavior, and convergence conditions: SymPy issue 8412 reports a summation producing \(\operatorname{lerchphi}(1,2,a)\) followed by an endpoint substitution that gives \texttt{nan}, and issue 14219 includes finite sums containing \(\operatorname{lerchphi}(x,1,k+2)\) with convergence restrictions around \(|x|\leq 1\) and \(x\neq 1\). Those reports are related to the broader LerchPhi endpoint/convergence family, but the artifact's best-effort deduplication check did not identify a public report for the direct left-hand limit \(\operatorname{limit}(\operatorname{lerchphi}(x,1,1),x,1,\mathrm{dir}=-)\) returning the unevaluated divergent endpoint. This exact reproducible case remains individually actionable as a SymPy issue and regression test.

\subsection{Proposed SymPy Regression Test}

\medskip

\begin{lstlisting}[style=bugpython]
import pytest

from sympy import S, lerchphi, limit, symbols


@pytest.mark.parametrize("a", [1, 2, 3, 4, 5, 6])
def test_left_limit_lerchphi_s_one_at_one_is_positive_infinity(a):
    x = symbols("x")
    assert limit(lerchphi(x, 1, a), x, 1, dir="-") == S.Infinity
\end{lstlisting}

\end{document}
